% ============================================================
%  R. Nielsen, "Lectures on 'Philosophical Propaedeutics'"
%  from the Academic Year 1860–61
%  Gyldendalske Boghandling, Copenhagen, 1862
%
%  Excerpt: pp. 326–328
%  Translation and introductory note by H. Halvorson, 2026.
%
%  Translator's note on terminology: Nielsen's technical
%  vocabulary is rendered as consistently as possible.
%  Paradigme: "paradigm" (his term for the organic archetype or
%  form-principle); Anlæg: "disposition"; Selvhed: "selfhood";
%  Idee: "Idea" (capitalized, following his Hegelian usage);
%  Andethed: "otherness"; Concretionsloven: "the law of
%  concretion." The German passages — Nielsen's quotations from
%  Bronn's translation of Darwin — are rendered into English
%  directly rather than back-translated through the Danish.
% ============================================================
\documentclass[12pt,a4paper]{article}
\usepackage[utf8]{inputenc}
\usepackage[T1]{fontenc}
\usepackage[english]{babel}
\usepackage{microtype}
\usepackage{setspace}
\usepackage[margin=1.2in]{geometry}
\usepackage{fancyhdr}
\setlength{\headheight}{14.5pt}
\usepackage{hyperref}

\hypersetup{
  pdftitle={Lectures on Philosophical Propaedeutics (excerpt)},
  pdfauthor={R. Nielsen},
  hidelinks
}

\pagestyle{fancy}
\fancyhf{}
\rhead{R.\ Nielsen, \emph{Phil.\ Prop.}\ (1862), pp.\ 326--328 [trans.\ HH]}
\cfoot{\thepage}

\setstretch{1.3}

\title{\textbf{Lectures on ``Philosophical Propaedeutics''}\\[0.4em]
\large from the Academic Year 1860--61\\[0.3em]
\normalsize Excerpt: pp.\ 326--328\\
(on the origin of species and Darwin's theory)\\[0.6em]
\small Translation and note by H.\ Halvorson (2026)}
\author{R.\ Nielsen}
\date{Gyldendalske Boghandling (F.\ Hegel), Copenhagen, 1862}

\begin{document}
\maketitle
\thispagestyle{fancy}

% ------------------------------------------------------------------
%  CONTEXTUALIZING NOTE
% ------------------------------------------------------------------

\section*{Note on Context}

Rasmus Nielsen (1809--1884) held the chair in philosophy at the
University of Copenhagen from 1850 until his death, and these
lectures from the academic year 1860--61 represent his most
systematic attempt to define the relationship between philosophy
and the natural sciences. He knew Darwin's \emph{Origin of
Species} (1859) through H.\ G.\ Bronn's German translation
(\emph{Ueber die Entstehung der Arten im Thier- und
Pflanzenreich}, Stuttgart 1860), which appeared just one year
after the English original and nearly a decade before J.\ P.\
Jacobsen's Danish translation (1871--72). The passage translated
here is, to our knowledge, the earliest substantive philosophical
engagement with Darwin in the Danish tradition.

The excerpt falls at the end of the longest section of the book,
§II.C, ``The Relation of Philosophy to Physiology: The Problem of
Life.'' To understand what Nielsen is doing with Darwin, one needs
to know where he is in the argument.

\medskip\noindent\textbf{The problem of organic types.} Nielsen's
central concern throughout §II.C is to establish that organic
types --- species, genera, families --- have genuine reality,
against two opposite errors. The first error is crude empiricism,
represented by the botanist Matthias Schleiden: only individuals
exist; species and genera are mental abstractions we project onto
nature. The second error is dogmatic Naturphilosophie: the type
exists as a free-standing metaphysical entity (an
\emph{Urwesen}, a primordial being) that imposes form on
matter from outside. Nielsen's alternative, developed at length in
the preceding pages, is that organic types have \emph{dialectical}
reality. The species-idea (\emph{Artsidee}) is not an abstraction
from individuals but their causal ground --- more concrete, not
less, than any individual, because it is the richer principle in
which all individual differences are pre-formed and, after their
development, remembered.

\medskip\noindent\textbf{The law of concretion.} This culminates
in what Nielsen calls \emph{Concretionsloven}, the law of
concretion: genus is more concrete than species, family more
concrete than genus, and so on up to the ground-idea of the
kingdom (\emph{Rigets Grundidee}). Crucially, however, this
hierarchy cannot be derived a priori. There is no philosophical
law determining how many species must belong to each genus, or
which genera form a family. Philosophical cognition moves
upward from empirical facts to concrete ideas, not downward
from abstract principles to facts. This anti-deductivist
commitment means that for Nielsen, biology genuinely constrains
philosophy --- he is not doing speculative Naturphilosophie in
the Schelling--Oken mold.

\medskip\noindent\textbf{Darwin's entry.} It is at this point ---
with the positive argument complete --- that Nielsen introduces
Darwin, and the move is philosophically precise. He does not
reject Darwin's empirical findings; on the contrary, he thinks
the observed interaction between organic type and environmental
conditions is exactly what his own dialectical theory predicts.
What he resists is a particular \emph{interpretation} of Darwin:
the reading that external conditions \emph{constitute} organic
types, that environment is all the way down. Such a reading would
make the type a mere resultant of circumstances, which for Nielsen
dissolves the very concept of life --- life without an ideal causal
principle is not life but chemistry.

The conditional framing of his objection (``if... the attempts...
are to be understood as meaning that it is actually the
circumstances and external conditions that determine the organic
types'') is therefore deliberate. Nielsen leaves open whether
Darwin himself holds this view. His target is a philosophical
conclusion that might be drawn from Darwin, not Darwin's science
as such.

\medskip\noindent\textbf{Bronn's \emph{Schlußwort}.} The long
note that follows in the original --- so extensive that it fills
an entire printed page --- quotes from H.\ G.\ Bronn's critical
afterword to his German translation. Bronn was a distinguished
paleontologist who accepted common descent but was skeptical of
natural selection as a complete explanation. His objection ---
that natural selection cannot explain \emph{why} environmental
pressure produces oval rather than lanceolate leaves, or five
rather than four stamens --- is, from Nielsen's perspective, the
right objection, though made for the wrong reasons. Bronn is
pointing to the absence of any intelligible connection between
the external condition and the specific organic result. For
Nielsen, what is missing is precisely the ideal causal principle,
the paradigm: without it, the specificity of organic form
remains unintelligible, a brute fact. Natural selection can
select among variations, but it cannot explain the determinate
character of those variations, and that explanatory gap is where
Nielsen's philosophy of the organism is meant to fit.

\medskip\noindent\textbf{The turn to p.\ 328.} After the Darwin
interlude, Nielsen pivots to a parallel critique of the opposite
error. The passage on pp.\ 327--328 attacks both Kantian critical
philosophy (which makes the Idea a merely subjective, regulative
principle) and dogmatic positivism (which hypostatizes it into a
free-standing primordial being). Darwin's error, on the
``lawless'' reading, belongs to the positivist side: dissolving
type into external conditions, making the observable phenotype
the only reality. The Naturphilosoph's error is the mirror image:
fixing type as a dogmatic immediate, unmediated by the critical
reflection that alone gives it scientific significance. Nielsen's
dialectical Idealism --- shaped, he believed, by Kierkegaard's
critique of system-building --- is meant to hold both sides in
tension without collapsing into either.

\medskip\noindent\textbf{Significance.} This passage shows that
engagement with Darwin was already philosophically serious in
Copenhagen in the early 1860s, more than a decade before the
Brandes generation made Darwin a cultural-political weapon. For
Nielsen, Darwin poses not a crisis of faith but a philosophical
problem about the intelligibility of life --- a problem that
requires more careful thinking about causality, not the
abandonment of teleology. His position anticipates, in some
respects, later neo-Kantian and vitalist responses to Darwinism
(Driesch, Bergson), though Nielsen's framework is more directly
Hegelian. It is also worth noting that Nielsen's student
Harald H\o ffding would work through many of the same questions
in ``Filosofien og Darwinismen'' (1874), the primary reading for
this week --- but from a much more empiricist and
epistemologically cautious vantage point.

\bigskip
\hrule
\bigskip

% ------------------------------------------------------------------
%  TRANSLATION
% ------------------------------------------------------------------

\section*{Translation}

\noindent\textit{[P.\ 326, continued from p.\ 325:]}

\medskip

\noindent through all its species to the family, etc., can be
determined dialectically by the law of concretion; on the other
hand, there is no law of development by which it could be
determined a priori how many and which species must necessarily
belong to each genus, or how many and which genera originally
belong to each family: philosophical cognition does not proceed
through abstract Ideas to actual facts, but conversely through
actual facts to the concrete Ideas. (Cf.\ \emph{Phil.\ Prop.},
pp.\ 158--178.)

If the attempts that have been made in recent times to explain the
origin of species from varieties that have become constant over
time are to be understood as meaning that it is actually the
circumstances and the external conditions that determine the
organic types, then such a lawless and Idea-negating conception of
vital activity will naturally meet with objection in philosophy.
But it ought not to be forgotten in this connection that the
reality of Ideas is infinitely dialectical, and that a dialectical
reality does not admit of being hypostatized in dogmatic positive
immediacy. If it is really the case that organic selfhood
originally presupposes elementary otherness, then it follows
directly from this that the development of life must also bring
with it a manifold interaction of the typical and the elementary.
Varieties, abortions, displacements, and other abnormalities show
clearly enough what influence the elementary conditions must have
on the execution of the paradigm. [*]

\medskip
\noindent\textit{[*] Note (pp.\ 326--327 in the original): For
the investigations pertaining to this matter, we refer to Charles
Darwin, \emph{On the Origin of Species in the Animal and Plant
Kingdom}, translated from the English and annotated by Dr.\ H.\
G.\ Bronn (Stuttgart, 1860). The following is quoted from the
``Concluding Remarks of the Translator,'' pp.\ 508--9:}

\begin{quote}
\itshape
``These are the conditions of existence to which organisms must
adapt themselves, and which play such a large role in this book.
They are partly organic and partly inorganic, and the former are,
in Darwin's view, far the most numerous and important, and for
that reason also in themselves most capable of giving rise to the
most varied consequences. But these very organic principles have
in turn the great advantage for Darwin that, since he appeals to
their variety and to the struggle for existence in general, he is
relieved of the necessity of accounting for their mode of
operation in detail and of showing what specific consequences this
or that specific organic condition exerts on the structure and
development of the organisms subject to its influence, either in
general or individually. Mr.\ Darwin appeals on every page to
the fact that only those variations have any prospect of survival
that are beneficial to the individual and thus to the future
species; and theoretically one must concede that, if natural
selection exists, the matter cannot behave otherwise. But we must
confess that in almost all of our varieties said to have arisen
from inner causes we cannot find at all wherein the benefit of
their modification consists\ldots. Why does one plant species
acquire, in this struggle for existence, oval instead of
lanceolate leaves, and another lanceolate instead of oval? Why
does one have an umbel-like and another a paniculate
inflorescence? Why does one have five and another four stamens,
one a closed and another a wide-open flower? Of what use to one
is this, and to the other the opposite? Why do the organic
conditions bring about this? By what means do they set about it,
and how must they be constituted to be capable of doing so? And
how can the one species thereby become superior to the other? We
confess that we cannot discern any connection between these
phenomena, and Mr.\ Darwin would answer us that it might possibly
happen in one way or another. We confess further that we have
looked in vain for positive proofs or even supporting evidence in
this regard, with the exception perhaps of some special cases of
a peculiar kind; for we are far from wishing to deny all such
influence altogether.''
\end{quote}

\noindent\textit{[P.\ 327: The body text begins at the top of
the page with three lines immediately following the end of the
note:]}

\medskip

\noindent ``The Critique of Pure Reason'' transforms the Idea
into a subjectively unconditioned entity; dogmatic positivism
hypostatizes the Idea into a primordial being subsisting in
itself. Such a positive

\medskip
\noindent\textit{[P.\ 328, continued from p.\ 327:]}

\medskip

\noindent primordial being is, to be sure, declared to be an
inward, imageless, and invisible being, but becomes, precisely
when its positivity is fixed, through its dogmatic immediacy an
outward, pictorial, and perceptible being. Deep Intuition, the
immediate Intuition of Ideas, is in general a mystical-dogmatic
intuition of a Negative-Positive, an Imageless-Pictorial; the
Depth is recognized above all by the fact that the one who
intuits does not notice the contradiction. If the immediate
Intuition of Ideas is to acquire scientific significance, this
can only come about through Critique not merely pointing out but
also thoroughly reflecting the contradiction that breaks through
from all sides. The contradiction is that the Immediate in
Nature is so reflected; the contradiction is that what is thus
Reflected is so immediate. Just as exact infinitesimal analysis
is not finished with the insight that every finite magnitude is
conceived as an integral of infinitely small magnitudes, just so
little is the dialectical analysis of Ideas concluded with the
insight that the Immediate reflects itself in the Reflected; for
the articulating intermediate determinations that develop from
this are innumerable.

The realm of Ideas is, in relation to the external plant-world
grounded therein, a system of reflection-existences, content-rich
existence-memories; the reflection of Ideas in the reflected
immediacy of individual forms is, like the clarity of light, a
transfigured immediacy. In this clarity the individual existences
are seen, through their individual Ideas, reflected in the Idea
of the species; the Ideas of species, with their individual
Ideas, reflected in the Idea of the genus; the Ideas of genera,
with their Ideas of species, in that of the family; through the
family's in the order's; through the order's in the class's;
through the class's in the series'; through the series' concrete,
content-rich Idea in the Idea that grounds the kingdom, in the
Depth.

\medskip
\noindent $\beta$) \textbf{Animal Life.}

As organic Nature, the Idea is Life, and living Existence,
essentially viewed, an Idea-Existence; but Idea-Existence

\medskip
\noindent\textit{[Continued on p.\ 329, not included here.]}

\end{document}
